\section{Related Work}
\label{sec:related_work}

\paragraph{VLM Hallucination.}
Hallucination in VLMs, where generated text is factually inconsistent with the visual input, has been studied along several mitigation strategies~\cite{li2023pope,wang2023amber}.
Supervised alignment methods adapt RLHF or DPO with human-annotated preference data to reduce hallucinated outputs~\cite{sun2023llava_rlhf,yu2023rlhfv,zhao2024hdpo,li2025dpo,ying2025empo}, but their reliance on costly labels limits scalability to large unlabeled corpora.
Training-free inference-time approaches, such as contrastive decoding~\cite{leng2023vcd}, attention-based penalty strategies~\cite{huang2023opera}, and dynamic correction decoding~\cite{wang2024deco,yin2025clearsight,tang2025farsight,anand2026crops}, intervene at generation without modifying model weights, yet leave underlying hallucination tendencies unchanged.
Representation-editing methods~\cite{jiang2024interp,wang2026after} steer internal activations toward factual semantics with low overhead, but require explicit factual supervision to define the editing direction.
% Closer to our setting, annotation-free self-training approaches such as STIC~\cite{deng2024stic} construct preference data from corrupted-image negatives, and APASI~\cite{lu2025apasi} self-injects hallucinations to form training pairs; visual evidence prompting~\cite{li2025visual_evidence} further grounds generation in salient image regions.
Unlike these methods, \ours constructs fine-grained rewards entirely from the model's own outputs via a regulated self-verification pipeline, requiring no human annotation, no external verifier, and no corrupted-image heuristics, while the VCR mechanism further ensures reward reliability against redundant and image-agnostic claims.


%% 阅读下面的文章，分类，并根据我的文章的内容，筛选润色，写一段related work，最后突出我的方法的创新点，和现有方法的区别
% https://arxiv.org/abs/1809.02156
% https://arxiv.org/abs/2305.10355
% https://arxiv.org/abs/2311.07397
% https://arxiv.org/abs/2309.14525 Aligning Large Multimodal Models with Factually Augmented RLHF
% https://arxiv.org/abs/2310.16045
% https://arxiv.org/abs/2311.17911
% https://arxiv.org/abs/2311.16922
% https://arxiv.org/abs/2410.11779
% https://arxiv.org/abs/2505.16652
% https://arxiv.org/abs/2601.00659
% https://aclanthology.org/2025.acl-long.634/
% https://arxiv.org/abs/2601.01957
% https://arxiv.org/abs/2505.16411
% https://arxiv.org/abs/2503.13107
% https://arxiv.org/abs/2411.10436
% https://arxiv.org/abs/2501.09695
% https://arxiv.org/abs/2506.04039
% https://arxiv.org/abs/2405.19716 STIC
% https://arxiv.org/abs/2509.11287
% https://arxiv.org/abs/2605.00323


% \subsection{Hallucination in VLMs}
% Hallucination in Vision-Language Models (VLMs) refers to the misalignment between 
% the factual content of images and the textual output generated by the model~\cite{liu2024survey}. 
% Such hallucinations manifest across multiple dimensions, including object existence, 
% attribute descriptions, and inter-object relations. For instance, a model may fabricate 
% non-existent objects, misidentify colors or quantities, or incorrectly describe spatial 
% relationships. These issues arise from a combination of factors, including biased training 
% data, insufficient visual grounding, and cross-modal misalignment~\cite{liu2024survey}. 
% Existing mitigation strategies broadly fall into three categories: training data 
% optimization, model architecture refinement, and post-processing of generated outputs. 
% However, most approaches rely on human-annotated labels or strong external supervision, 
% limiting their scalability to large unannotated datasets. In contrast, our method 
% constructs a fully self-supervised reward signal directly from the model's visual and 
% textual outputs, requiring no ground-truth annotation.

% \subsection{Annotation-Free Training for VLMs}
% Recent work has explored training VLMs without human-annotated labels by constructing 
% automatic supervision signals. Vision-Zero~\cite{wang2026visionzero} proposes a 
% label-free, multi-agent self-play framework where VLMs engage in a ``Who Is the Spy?'' 
% visual game to generate their own training data. By alternating between self-play and 
% RLVR via their Iterative-SPO algorithm, Vision-Zero achieves sustained performance 
% gains on reasoning and vision-centric tasks without any human annotation. While 
% Vision-Zero focuses on constructing competitive game environments to elicit general 
% reasoning abilities, our work takes a different approach: we design a fact-checking 
% pipeline that directly targets visual hallucination by generating verifiable reward 
% signals from the model's own outputs. Rather than relying on inter-agent competition, 
% FactReward grounds the reward in factual consistency between generated descriptions 
% and the original image, providing a more targeted training signal for reducing image hallucinations.
